\section{Experimental Design}

\subsection{Datasets, Detection Methods, and Evaluation Metrics}

\subsubsection{POPE: Object Existence Hallucination Detection}

The detection method for POPE is rule-based matching.
During response generation,
the model receives a prompt containing only the question
and an instruction to answer Yes or No,
ensuring concise and parseable responses
(the complete prompt template is provided in \Cref{sec:appendix}).
The answer post-processing algorithm follows
the official evaluation script from the
RUCAIBox/POPE repository;\footnote{\url{https://github.com/RUCAIBox/POPE/blob/main/evaluate.py}}
its pseudocode is given in \Cref{alg:pope_normalize}.

\begin{algorithm}[htbp]
  \caption{POPE answer normalization}
  \label{alg:pope_normalize}
  \begin{algorithmic}[1]
    \Procedure{NormalizePOPEAnswer}{response}
      \State sentence $\gets$ first sentence of response (split by ``.'')
      \State remove all commas from sentence
      \State words $\gets$ split sentence by spaces
      \If{``No'' $\in$ words \textbf{or} ``not'' $\in$ words \textbf{or} ``no'' $\in$ words}
        \State \Return ``no''
      \EndIf
      \State \Return ``yes''
    \EndProcedure
  \end{algorithmic}
\end{algorithm}

Hallucination is defined following the POPE protocol:
a sample is counted as an object hallucination (false positive)
when the model answers ``yes'' but the ground truth label is ``no,''
corresponding to object-level faithfulness hallucination.

Evaluation metrics treat ``yes'' (object present) as the positive class,
so object hallucination corresponds to false positives.
The core metrics are summarized in \Cref{tab:pope_metrics}.

\begin{table}[htbp]
  \centering
  \caption{POPE evaluation metrics.}
  \label{tab:pope_metrics}
  \begin{tabular}{lll}
    \toprule
    Metric & Formula & Description \\
    \midrule
    Accuracy          & $(TP+TN)/N$        & Overall correctness \\
    Precision         & $TP/(TP+FP)$       & Accuracy when predicting Yes \\
    Recall            & $TP/(TP+FN)$       & Detection rate for actual Yes \\
    F1                & $2PR/(P+R)$        & Harmonic mean of Precision and Recall \\
    Yes Ratio         & $N_{\text{yes}}/N$ & Response bias indicator \\
    Hallucination Rate & $FP/N$            & Core metric; equals object hallucination rate \\
    \bottomrule
  \end{tabular}
\end{table}

\subsubsection{MathVista: Mathematical Reasoning Hallucination Detection}

The detection method for MathVista employs an MLLM Judge
following the MMHal-Bench evaluation protocol
\citep{sunAligningLargeMultimodal2023},
which adopts a 0--6 Likert scale that jointly assesses
the informativeness and hallucination level of each response,
using a threshold of 3 (scores below 3 indicate hallucination).
In our implementation, GPT-5.5 serves as the judge model.
The judge's system prompt defines the scoring rubric,
hallucination criteria,
and a three-way hallucination type classification,
accompanied by five few-shot examples to ensure consistency.
When the judge assigns a score below 3,
it additionally labels the hallucination type
(faithfulness, factuality, or logical).
The complete prompt templates for both model response generation
and judge scoring are provided in \Cref{sec:appendix}.

During response generation,
prompts are dynamically constructed from sample metadata
to produce format-specific instructions.
In Direct mode,
the prompt consists of the question text
concatenated with format instructions.
In CoT mode,
the prompt further requires step-by-step reasoning
with the final answer marked by ``Final answer:''.
Both modes are evaluated
to analyze the effect of explicit reasoning on hallucination behavior.

The evaluation metrics are shown in \Cref{tab:mathvista_metrics}.

\begin{table}[htbp]
  \centering
  \caption{MathVista evaluation metrics.}
  \label{tab:mathvista_metrics}
  \begin{tabular}{lll}
    \toprule
    Metric & Definition & Description \\
    \midrule
    Hallucination Rate & $N_{\text{score}<3} / N_{\text{total}}$ & Lower is better \\
    Average Score & $\sum \text{score}_i / N$ & Overall response quality \\
    Type-specific count & By faithfulness/factuality/logical & Hallucination source breakdown \\
    \bottomrule
  \end{tabular}
\end{table}

\subsubsection{VQA-RAD: Medical Visual Question Answering Hallucination Detection}

The detection method for VQA-RAD follows the same MLLM Judge protocol
(0--6 Likert scale, threshold 3)
with domain-specific prompt adaptations for medical scenarios
(see \Cref{sec:appendix}).
The judge's task description is replaced
with a radiology-specific scenario description
that requires the judge to verify consistency
with both the image's visual content
and established medical and anatomical knowledge.
During response generation,
a radiologist expert role prompt is adopted
following the role-setting strategy
of \citet{panditHALTMedVQAHallucinationAware2025}.
The three-way hallucination type classification
(faithfulness, factuality, logical) is always enabled,
and hallucination rates are additionally stratified
by closed-ended and open-ended answer types.

The evaluation metrics include Hallucination Rate, Average Score,
type-specific hallucination counts,
and stratified metrics for closed-ended and open-ended questions.

\subsection{Experimental Conditions}

All four models are accessed through APIs
with a unified random seed of 42.
Temperature is kept at each platform's default,
and maximum token length is uncapped.
POPE samples 1{,}000 items per split,
MathVista uses the full testmini set (approximately 1{,}000 items),
and VQA-RAD uses the test split (451 items).
The MLLM Judge is GPT-5.5.

\Cref{tab:method_dataset_overview} summarizes
the correspondence among detection methods, datasets, and metrics.

\begin{table}[htbp]
  \centering
  \caption{Overview of method--dataset--metric correspondence.}
  \label{tab:method_dataset_overview}
  \begin{tabular}{llll}
    \toprule
    Dataset   & Detection Method & Hallucination Types  & Core Metrics \\
    \midrule
    POPE      & Rule-based       & Faithfulness         & HR, F1 \\
    MathVista & MLLM Judge       & All three types      & HR, Avg Score \\
    VQA-RAD   & MLLM Judge       & All three types      & HR, Avg Score, Closed/Open HR \\
    Human     & Manual annotation & All three types     & Kappa, Pearson $r$, F1 \\
    \bottomrule
  \end{tabular}
\end{table}
